%% Die Sprachdatei von helsinki für en-US (ADR-0025).
%% Jedes Wort wie en-GB, aber eine eigene Datei ohne Rückfall, damit eine fehlende Übersetzung auffällt.

%% ---------------------------------------------------------------- Zahlen
%% \betrag wie in helsinki-en-gb.tex: €8,337.50 und −€200.00.
\sisetup{
  output-decimal-marker = {.},
  group-separator       = {,},
}
\newcommand{\zeichenEUR}{€}
\newcommand{\zeichenUSD}{\$}
\newcommand{\zeichenGBP}{£}
\newcommand{\waehrungszeichen}{\csname zeichen\waehrung\endcsname}
\makeatletter
\newcommand{\betrag}[1]{\betrag@zeichen#1\relax}
\def\betrag@zeichen#1#2\relax{\ifx-#1\textminus\waehrungszeichen\num{#2}\else\waehrungszeichen\num{#1#2}\fi}
\makeatother
\newcommand{\eurobetrag}[1]{\begingroup\def\waehrung{EUR}\betrag{#1}\endgroup}
\newcommand{\prozent}[1]{\num{#1}\%}

%% ------------------------------------------------------------- Seiten
\newcommand{\wortseitevon}[2]{Page #1 of #2}
\newcommand{\wortfortsetzung}{Continued on the next page}

%% ------------------------------------------------------ Kopf und Anschrift
\newcommand{\wortgeschaeftsfuehrung}{Managing director}
\newcommand{\wortustid}{VAT ID}
\newcommand{\worttelefon}{Tel.}
\newcommand{\wortkundennr}{Customer no.:}
\newcommand{\wortauftragsnr}{Our order no.:}
\newcommand{\wortvereinnahmtam}{Payment received on}
\newcommand{\wortreferenz}{Your reference:}
\newcommand{\wortzahlbarbis}{Due in full by}

%% ----------------------------------------------------- Positionen, Summen
\newcommand{\wortbeschreibung}{Description}
\newcommand{\wortmenge}{Quantity}
\newcommand{\worteinzelpreis}{Unit price}
\newcommand{\wortrabatt}{Discount}
\newcommand{\wortbetrag}{Amount}
\newcommand{\wortnettobetrag}{Net amount}
%% Satz und, bei mehreren Steuergruppen, ihre Grundlage – sonst ist sie leer.
\newcommand{\wortsteuer}[2]{plus #1 VAT\ifblank{#2}{}{ on #2}}
\newcommand{\wortgesamtbetrag}{Total amount}

%% --------------------------------------------------------- Schluss und Fuß
\newcommand{\wortgruss}{Kind regards}
\newcommand{\wortbankverbindung}{Bank details}
\newcommand{\wortkontoinhaber}{Account holder}
\newcommand{\wortinstitut}{Bank}
\newcommand{\wortregistergericht}{Registry court}
